\begin{lemma}
\label{lemma-henselization-nil}
Let $R$ be a local ring. Let $nil(R)$ denote the ideal of
nilpotent elements of $R$. Then $nil(R)R^h = nil(R^h)$ and
$nil(R)R^{sh} = nil(R^{sh})$.
\end{lemma}

\begin{proof}
Note that $nil(R)$ is the biggest ideal consisting of nilpotent elements
such that the quotient $R/nil(R)$ is reduced. Note that $nil(R)R^h$
consists of nilpotent elements by
Algebra, Lemma \ref{algebra-lemma-locally-nilpotent}.
Also, note that $R^h/nil(R) R^h$ is the
henselization of $R/nil(R)$ by
Algebra, Lemma \ref{algebra-lemma-quotient-henselization}.
Hence $R^h/nil(R)R^h$ is reduced by
Lemma \ref{lemma-henselization-reduced}.
We conclude that $nil(R) R^h = nil(R^h)$ as desired.
Similarly for the strict henselization but using
Algebra, Lemma \ref{algebra-lemma-quotient-strict-henselization}.
\end{proof}
